\documentclass[tikz,border=8pt]{standalone}
\usepackage[no-math]{fontspec}
\usepackage[brazil]{babel}
\usepackage{amsmath,amssymb}
\IfFileExists{C:/Windows/Fonts/Lato_Regular.ttf}{%
  \setmainfont{Lato}[
    Path = C:/Windows/Fonts/,
    UprightFont = Lato_Regular.ttf,
    BoldFont = Lato_Bold.ttf,
    ItalicFont = Lato_Italic.ttf,
    BoldItalicFont = Lato_BoldItalic.ttf
  ]
  \setsansfont{Lato}[
    Path = C:/Windows/Fonts/,
    UprightFont = Lato_Regular.ttf,
    BoldFont = Lato_Bold.ttf,
    ItalicFont = Lato_Italic.ttf,
    BoldItalicFont = Lato_BoldItalic.ttf
  ]
}{%
  \RequirePackage[default]{lato}
}
\renewcommand{\familydefault}{\sfdefault}

\usepackage{xcolor}
\usepackage{tikz}
\usepackage{pgfplots}
\pgfplotsset{compat=1.18}
\usetikzlibrary{arrows.meta,positioning,shapes.geometric,fit,backgrounds,calc,decorations.pathreplacing,patterns}

% Paleta institucional UEFS / CHF662
\definecolor{uefsNavy}{HTML}{1B2A4A}
\definecolor{uefsTeal}{HTML}{0E7C7B}
\definecolor{uefsForest}{HTML}{2D6A4F}
\definecolor{uefsGold}{HTML}{C4983B}
\definecolor{uefsGoldDark}{HTML}{8F6B1E}
\definecolor{uefsAlert}{HTML}{B02A37}
\definecolor{uefsOrange}{HTML}{D95F02}
\definecolor{uefsInk}{HTML}{22252A}
\definecolor{uefsMuted}{HTML}{5C6470}
\definecolor{uefsWater}{HTML}{2E86AB}
\definecolor{uefsSand}{HTML}{D9B382}
\definecolor{bgNavy}{HTML}{F2F5F9}
\definecolor{bgTeal}{HTML}{E9F5F4}
\definecolor{bgForest}{HTML}{EEF5F0}
\definecolor{bgGold}{HTML}{FCF7EB}
\definecolor{bgAlert}{HTML}{FDF2F2}
\definecolor{bgGray}{HTML}{F8FAFC}
\definecolor{lineGray}{HTML}{CFD8DC}

\tikzset{
  >=Stealth,
  aaCard/.style={draw=#1, fill=white, rounded corners=4pt, line width=1pt, inner sep=6pt, align=center, font=\small},
  aaBox/.style={draw=#1, rounded corners=4pt, line width=1.1pt, inner sep=6pt, align=center, font=\footnotesize\sffamily, text=uefsInk},
  aaTitle/.style={font=\normalsize\bfseries\sffamily, text=uefsNavy, align=center},
  aaSub/.style={font=\footnotesize\sffamily, text=uefsMuted, align=center},
  aaArrow/.style={->, >=Stealth, line width=1.4pt},
  aaDArrow/.style={<->, >=Stealth, line width=1.4pt}
}


\begin{document}
\begin{tikzpicture}[font=\sffamily]

% Título principal
\node[font=\large\bfseries\sffamily, text=uefsNavy, align=center] at (0, 4.25)
  {Sistema Brasileiro de Licenciamento Ambiental (Trifásico \& Modalidades)};

% =========================================================
% PIPELINE SUPERIOR: LP -> LI -> LO
% =========================================================
% LP
\begin{scope}[shift={(-5.7, 1.95)}]
  \draw[draw=uefsNavy, fill=bgNavy, rounded corners=6pt, line width=1.2pt] (-2.4, -1.45) rectangle (2.4, 1.55);
  \fill[uefsNavy, rounded corners=5pt] (-2.4, 0.9) rectangle (2.4, 1.55);
  \fill[uefsNavy] (-2.4, 0.9) rectangle (2.4, 1.15);
  \node[font=\small\bfseries\sffamily, text=white] at (0, 1.22) {LP — Licença Prévia};

  \node[anchor=north west, text width=4.35cm, align=left, font=\footnotesize\sffamily, text=uefsInk] at (-2.2, 0.72) {%
    $\bullet$ \textbf{Viabilidade ambiental} e\\[-1pt]\hspace{0.55em}locacional na concepção\\[4pt]
    $\bullet$ Aprova estudos (\textbf{EIA/RIMA})\\[4pt]
    $\bullet$ Estabelece condicionantes
  };
\end{scope}

% LI
\begin{scope}[shift={(0, 1.95)}]
  \draw[draw=uefsGoldDark, fill=bgGold, rounded corners=6pt, line width=1.2pt] (-2.4, -1.45) rectangle (2.4, 1.55);
  \fill[uefsGoldDark, rounded corners=5pt] (-2.4, 0.9) rectangle (2.4, 1.55);
  \fill[uefsGoldDark] (-2.4, 0.9) rectangle (2.4, 1.15);
  \node[font=\small\bfseries\sffamily, text=white] at (0, 1.22) {LI — Licença de Instalação};

  \node[anchor=north west, text width=4.35cm, align=left, font=\footnotesize\sffamily, text=uefsInk] at (-2.2, 0.72) {%
    $\bullet$ Autoriza \textbf{início de obras}\\[-1pt]\hspace{0.55em}e implantação física\\[4pt]
    $\bullet$ Aprova planos e \textbf{PBA}\\[4pt]
    $\bullet$ Verifica requisitos da LP
  };
\end{scope}

% LO
\begin{scope}[shift={(5.7, 1.95)}]
  \draw[draw=uefsForest, fill=bgForest, rounded corners=6pt, line width=1.2pt] (-2.4, -1.45) rectangle (2.4, 1.55);
  \fill[uefsForest, rounded corners=5pt] (-2.4, 0.9) rectangle (2.4, 1.55);
  \fill[uefsForest] (-2.4, 0.9) rectangle (2.4, 1.15);
  \node[font=\small\bfseries\sffamily, text=white] at (0, 1.22) {LO — Licença de Operação};

  \node[anchor=north west, text width=4.35cm, align=left, font=\footnotesize\sffamily, text=uefsInk] at (-2.2, 0.72) {%
    $\bullet$ Autoriza \textbf{operação}\\[-1pt]\hspace{0.55em}efetiva do empreendimento\\[4pt]
    $\bullet$ Fiscaliza execução da LI\\[4pt]
    $\bullet$ Impõe \textbf{condicionantes} e\\[-1pt]\hspace{0.55em}monitoramento contínuo
  };
\end{scope}

% Setas do rito trifásico
\draw[->, >=Stealth, line width=2pt, uefsNavy] (-3.2, 1.95) -- (-2.5, 1.95);
\draw[->, >=Stealth, line width=2pt, uefsGoldDark] (2.5, 1.95) -- (3.2, 1.95);

% =========================================================
% DIVISOR INFERIOR: RITOS ESPECIAIS E COMPETÊNCIAS
% =========================================================
\draw[lineGray, line width=1pt] (-8.1, 0.05) -- (8.1, 0.05);
\node[fill=white, inner xsep=8pt, font=\footnotesize\bfseries\sffamily, text=uefsMuted] at (0, 0.05)
  {Ritos Especiais e Competências (SISNAMA)};

% Card (a): Modalidades
\node[draw=uefsTeal, fill=bgTeal, rounded corners=6pt, line width=1.2pt,
      anchor=north, text width=7.65cm, minimum height=2.85cm, inner sep=8pt, align=left,
      font=\footnotesize\sffamily, text=uefsInk] at (-4.25, -0.35)
  {\textbf{\small\color{uefsTeal}(a) Modalidades da Lei 15.190/2025 e Bahia}\\[2pt]
   \scriptsize\textbf{\color{uefsMuted}(Lei Estadual 10.431/2006 $\cdot$ INEMA/CEPRAM)}\\[8pt]
   $\bullet$ \textbf{LAC} (Adesão e Compromisso) e \textbf{LAU} (Única / Simplificada)\\[5pt]
   $\bullet$ \textbf{Licença Unificada (LU)} e \textbf{LOC} (Operação Corretiva)};

% Card (b): LC 140/2011
\node[draw=uefsNavy, fill=bgGray, rounded corners=6pt, line width=1.2pt,
      anchor=north, text width=7.65cm, minimum height=2.85cm, inner sep=8pt, align=left,
      font=\footnotesize\sffamily, text=uefsInk] at (4.25, -0.35)
  {\textbf{\small\color{uefsNavy}(b) LC 140/2011 (SISNAMA) — Competências}\\[8pt]
   $\bullet$ \textbf{Federal (IBAMA):} impacto nacional/regional ($\ge 2$ UFs)\\[5pt]
   $\bullet$ \textbf{Estadual (INEMA):} impacto estadual ou caráter supletivo\\[5pt]
   $\bullet$ \textbf{Municipal:} impacto ambiental estritamente local};

\end{tikzpicture}
\end{document}
